\documentclass[a4paper,10pt]{article}
\usepackage[czech]{babel}
\usepackage[T1]{fontenc}
\usepackage[utf8]{inputenc}
\usepackage{amsmath,amssymb}
\usepackage[margin=18mm]{geometry}
\setlength{\parindent}{0pt}
\begin{document}
\section*{Geometrická posloupnost C/D -- řešení a bodování}
Podklad pro učitele k testu
\texttt{01\_test\_GeometrickaPosloupnost\_CD.tex}.

Zařazení: 4. ročník SPŠ; tematický plán
\texttt{4\_rocnik/tematicke\_plany/SPS.pdf}, s.~1, září, bod 1.3;
celek posloupností má společnou dotaci 12 hodin.
Vzory z pracovního listu \texttt{pracovniListGeoPosl.tex}:
úlohy 6--7 (obecný člen), 24 (dvě znaménka kvocientu),
49 (kvocient ze součtu), 64--65 (součet a součin trojice).
V úloze 3 je tento vzor rozšířen o výpočet dalšího konečného součtu.
Nepotřebujeme logaritmy, goniometrii ani nekonečné řady.

Po vzoru A/B předpokládáme práci bez kalkulačky.
Odhad času je 30--35 minut včetně kontroly, nikoli původních 20 minut;
není ověřen ve třídě. Obě varianty mají shodný počet rozhodovacích kroků,
dvě řešení v úlohách 2--4 a stejné bodování.

\subsection*{Varianta C}
\begin{enumerate}
\item Pro každé celé $n\geq1$ je $4^n\ne0$ a $a_n\ne0$. Platí
\[
\frac{a_{n+1}}{a_n}
=\frac{3(-2)^{n+2}}{4^{n+1}}\frac{4^n}{3(-2)^{n+1}}
=-\frac12.
\]
Podíl je konstantní pro všechna $n\geq1$, takže posloupnost je geometrická,
$a_1=3(-2)^2/4=3$, $q=-\frac12$.
Ověření jen několika počátečních podílů by samo o sobě důkazem nebylo.

\item Z $a_4=a_2q^2$ plyne $3=12q^2$, tedy $q^2=\frac14$.
Člen $a_2$ je nenulový, proto $q\ne0$ a můžeme použít
$a_1=a_2/q$, $a_5=a_4q$:
\[
\begin{array}{c|c|c}
q&a_1&a_5\\ \hline
\frac12&24&\frac32\\
-\frac12&-24&-\frac32
\end{array}
\]
Rovnice pro $q^2$ má právě tyto dva reálné kořeny.
Zkouška postupným násobením:
\[
24,12,6,3,\tfrac32;\qquad -24,12,-6,3,-\tfrac32.
\]
Obě posloupnosti splňují zadání.

\item Součet prvních tří členů zapíšeme přímo, bez dělení výrazem $q-1$:
\[
2(1+q+q^2)=14
\iff q^2+q-6=0
\iff(q-2)(q+3)=0.
\]
Máme právě $q=2$ nebo $q=-3$; oba kvocienty se liší od jedničky.
Pro každou možnost zvlášť:
\[
\begin{aligned}
q=2:&\quad S_5=2\frac{2^5-1}{2-1}=62,\\
q=-3:&\quad S_5=2\frac{(-3)^5-1}{-3-1}=122.
\end{aligned}
\]
Zkouška přímým součtem:
\[
\begin{aligned}
2+4+8&=14,&2+4+8+16+32&=62,\\
2-6+18&=14,&2-6+18-54+162&=122.
\end{aligned}
\]

\item Označme trojici $u,v,w$. Součin je nenulový, takže všechny tři
členy i kvocient jsou nenulové. Z geometrické vlastnosti plyne $uw=v^2$,
a proto $uvw=v^3=216$. Jediné reálné řešení je $v=6$.
Zbývá $u+w=15$, $uw=36$. Dosazením $w=15-u$ dostaneme
\[
u(15-u)=36\iff u^2-15u+36=0\iff(u-3)(u-12)=0.
\]
Všechna řešení jsou $(3,6,12)$ s $q=2$ a $(12,6,3)$ s $q=\frac12$.
Zkouška: v obou případech je součet $21$, součin $216$ a oba sousední
podíly jsou stejné ($2$, resp. $\frac12$).
Jiné řešení není možné: $v$ je určen jednoznačně a rovnice pro $u$ má
právě dva kořeny. Zkouška prokazuje i postačitelnost získaných podmínek.
\end{enumerate}

\newpage
\subsection*{Varianta D}
\begin{enumerate}
\item Pro každé celé $n\geq1$ je $9^{n-1}\ne0$ a $a_n\ne0$. Platí
\[
\frac{a_{n+1}}{a_n}
=\frac{5(-3)^n}{9^n}\frac{9^{n-1}}{5(-3)^{n-1}}
=-\frac13.
\]
Podíl je konstantní pro všechna $n\geq1$, takže posloupnost je geometrická,
$a_1=5(-3)^0/9^0=5$, $q=-\frac13$.

\item Z $a_5=a_3q^2$ plyne $24=6q^2$, tedy $q^2=4$.
Člen $a_3$ je nenulový, proto $q\ne0$. Použijeme
$a_1=a_3/q^2$, $a_6=a_5q$:
\[
\begin{array}{c|c|c}
q&a_1&a_6\\ \hline
2&\frac32&48\\
-2&\frac32&-48
\end{array}
\]
Rovnice pro $q^2$ má právě tyto dva reálné kořeny.
Zkouška postupným násobením: $\frac32,3,6,12,24,48$,
resp. $\frac32,-3,6,-12,24,-48$. Obě posloupnosti splňují zadání.
Znaménko $q$ zde neovlivňuje $a_1$, ale ovlivňuje $a_6$.

\item Přímým zápisem součtu získáme
\[
3(1+q+q^2)=39
\iff q^2+q-12=0
\iff(q-3)(q+4)=0.
\]
Máme právě $q=3$ nebo $q=-4$; oba kvocienty se liší od jedničky.
\[
\begin{aligned}
q=3:&\quad S_5=3\frac{3^5-1}{3-1}=363,\\
q=-4:&\quad S_5=3\frac{(-4)^5-1}{-4-1}=615.
\end{aligned}
\]
Zkouška přímým součtem:
\[
\begin{aligned}
3+9+27&=39,&3+9+27+81+243&=363,\\
3-12+48&=39,&3-12+48-192+768&=615.
\end{aligned}
\]

\item Označme trojici $u,v,w$. Nenulový součin vylučuje nulové členy
i nulový kvocient. Platí $uw=v^2$, tedy $uvw=v^3=216$ a $v=6$.
Dále $u+w=20$, $uw=36$. Dosazením $w=20-u$ získáme
\[
u(20-u)=36\iff u^2-20u+36=0\iff(u-2)(u-18)=0.
\]
Všechna řešení jsou $(2,6,18)$ s $q=3$ a $(18,6,2)$ s $q=\frac13$.
Zkouška: v obou případech je součet $26$, součin $216$ a oba sousední
podíly jsou stejné ($3$, resp. $\frac13$).
Jiné řešení není možné: jediný reálný prostřední člen je $6$ a rovnice
pro krajní člen má právě dva kořeny. Obě trojice splňují původní zadání.
\end{enumerate}

\subsection*{Orientační bodování pro obě varianty}
Nezávazný návrh pro učitele; uznávejte i jiné správné postupy.
Kritéria lze upravit podle odborného posouzení. Stejnou početní chybu
netrestejte opakovaně, pokud navazující postup odpovídá získané hodnotě.
\begin{enumerate}
\item 3 body: obecný podíl a zdůvodnění, že je konstantní pro všechna
přípustná $n$ (1), první člen (1), kvocient (1).
\item 3 body: obě hodnoty kvocientu ze vztahu mezi známými členy (1),
odpovídající hodnoty $a_1$ (1), odpovídající hodnoty požadovaného členu (1).
\item 3 body: rovnice pro $S_3$ a obě hodnoty kvocientu (1), vzorec pro
$S_5$ se správným dosazením a úpravou mocnin (1), oba výsledky $S_5$ (1).
\item 4 body: geometrický vztah a určení prostředního členu (1),
obě uspořádané trojice (1), oba kvocienty (1), zkouška obou trojic (1).
\end{enumerate}
Celkem 13 bodů za každou variantu.
Orientační časy úloh: 4, 6, 8 a 9 minut, dále 3--8 minut na kontrolu.
Na dvacetiminutový test je tato skladba pravděpodobně příliš náročná.

\subsection*{Kontrola a rozdíl oproti A/B}
Výsledky byly nezávisle ověřeny přesnou racionální aritmetikou:
generováním členů z nalezených $a_1,q$, přímým sčítáním a dosazením
trojic do původních podmínek. V úloze 1 byla vedle obecného podílu
ověřena algebraická totožnost se vzorcem $a_1q^{n-1}$.
Úplnost řešení zajišťuje v úlohách 2 a 3 rozklad kvadratických rovnic,
v úloze 4 jediný reálný prostřední člen a oba kořeny rovnice pro krajní členy.
Žádným dělením se nevylučuje přípustné řešení; nenulovost je před dělením
zdůvodněna. V úloze 3 se podmínka $q\ne1$ používá až po určení kořenů.

Oproti A/B vyžaduje úloha 1 práci s mocninami a důkaz pro celý průběh
posloupnosti místo porovnání dvou číselných podílů.
V úloze 2 není kvocient zadán; žák ho odvodí ze dvou členů a musí
zpracovat obě znaménka. V úloze 3 nejprve řeší kvadratickou rovnici
pro kvocient a pak dva různé součty. V úloze 4 převádí současně součet,
součin a geometrickou podmínku na rovnici pro krajní členy;
na rozdíl od lineární rovnice v A/B musí zachytit obě pořadí trojice.
\end{document}
